% Created 2026-05-20 qua 22:36
% Intended LaTeX compiler: xelatex
\documentclass[12pt,a4paper]{article}
\usepackage{graphicx}
\usepackage{longtable}
\usepackage{wrapfig}
\usepackage{rotating}
\usepackage[normalem]{ulem}
\usepackage{capt-of}
\usepackage{hyperref}
\usepackage[T1]{fontenc}
\usepackage[utf8]{inputenc}
\usepackage[brazil, brazilian]{babel}
\usepackage[top=2.5cm,bottom=2.5cm,left=3cm,right=2.5cm]{geometry}
\usepackage{lmodern}
\usepackage{microtype}
\usepackage{minted}
\setminted{fontsize=\small, linenos=false, breaklines=true, bgcolor=gray!6}
\usepackage[colorlinks=true, linkcolor=black, urlcolor=blue, citecolor=black]{hyperref}
\usepackage{fancyhdr}
\pagestyle{fancy}
\fancyhf{}
\fancyhead[L]{\small Ciência de Dados - Laboratório 6}
\fancyhead[R]{\small Relatório de Modelagem}
\fancyfoot[C]{\thepage}
\renewcommand{\headrulewidth}{0.4pt}
\usepackage{setspace}
\setstretch{1.15}
\usepackage{fancyvrb}
\usepackage{etoolbox}
\usepackage{fancyvrb}
\DefineVerbatimEnvironment{Highlighting}{Verbatim}{commandchars=\\\{\}, formatcom=\small\color{gray}, frame=single, framesep=5mm, rulecolor=\color{lightgray}}
\usepackage{minted}
\usepackage{xcolor}
\definecolor{lightgray}{gray}{0.95}
\setminted{style=vs, fontsize=\small, frame=single, bgcolor=lightgray}
\AtBeginDocument{\setlength{\parskip}{4pt}}
\author{Nícolas Rosenthal Dal Corso}
\date{\today}
\title{Laboratório 6 - Relatório}
\hypersetup{
 pdfauthor={Nícolas Rosenthal Dal Corso},
 pdftitle={Laboratório 6 - Relatório},
 pdfkeywords={},
 pdfsubject={},
 pdfcreator={Emacs 30.2 (Org mode 9.7.11)}, 
 pdflang={Portuges}}
\begin{document}

\maketitle
\setcounter{tocdepth}{3}
\tableofcontents

\section{Introdução}
\label{sec:orgfc9dc0a}

Este relatório descreve o pipeline completo desenvolvido para o problema
\textbf{Spaceship Titanic}, da plataforma Kaggle.

O objetivo da competição consiste em prever se um passageiro foi
transportado para outra dimensão, representado pelo predicado

\begin{verbatim}
Transported ∈ {True, False}
\end{verbatim}

Partindo da análise exploratória de dados e da engenharia de \emph{features} desenvolvida anteriormente, o pipeline implementado utiliza \emph{ensemble stacking} e otimização Bayesiana com Optuna; validação cruzada estratificada, otimização de \emph{threshold}, e modelos \emph{gradient boosting} modernos (\texttt{CatBoost, LightGBM, XGBoost})

O foco principal do projeto foi maximizar desempenho preditivo, robustez estatística e capacidade de generalização.
\section{Estrutura Geral do Pipeline}
\label{sec:orgdd54f66}
O pipeline foi dividido em duas etapas principais:
\subsection{Pré-processamento em R}
\label{sec:org29d83c5}
\begin{itemize}
\item EDA (\emph{Exploratory Data Analysis});
\item imputação de valores ausentes;
\item transformação de variáveis;
\item feature engineering;
\item normalização estrutural;
\item exportação dos \emph{datasets} processados.
\end{itemize}

Arquivos gerados:

\begin{verbatim}
data/processed/train_processed.csv
data/processed/test_processed.csv
\end{verbatim}
\subsection{Modelagem em Python}
\label{sec:org1396ba2}
\begin{itemize}
\item \emph{encoding} categórico;
\item otimização de hiperparâmetros;
\item treinamento dos modelos;
\item \emph{stacking ensemble};
\item avaliação;
\item geração da submissão final.
\end{itemize}
\section{Pré-processamento e Engenharia de Features}
\label{sec:org09e5664}
\subsection{Imputação de Valores Ausentes}
\label{sec:org193da89}
Os valores ausentes foram tratados em R utilizando regras heurísticas
baseadas na análise exploratória.
\subsubsection{CryoSleep}
\label{sec:org432a181}
A \emph{feature} \texttt{CryoSleep} foi imputada utilizando a relação observada entre \textbf{ausência de gastos} e \textbf{estado criogênico}.
\subsubsection{VIP}
\label{sec:org9e8efda}
A variável \texttt{VIP} foi imputada utilizando distribuição empírica baseada em níveis elevados de gastos totais.
\subsubsection{Gastos}
\label{sec:orgdc81af9}
Passageiros em criossono tiveram os gastos redefinidos para zero:

\begin{itemize}
\item RoomService
\item FoodCourt
\item ShoppingMall
\item Spa
\item VRDeck
\end{itemize}
\subsection{Engenharia de Features}
\label{sec:orgc3fabbe}
\subsubsection{\texttt{Expenditure} e \texttt{NoExpenditure}}
\label{sec:org9fcffc5}

Soma total dos gastos do passageiro:
\begin{verbatim}
Expenditure =
 RoomService +
 FoodCourt +
 ShoppingMall +
 Spa +
 VRDeck
\end{verbatim}
\subsubsection{\texttt{Cabin} \emph{Features}}
\label{sec:org876baca}
A variável \texttt{Cabin} foi decomposta em \texttt{Cabin\_Deck}, \texttt{Cabin\_Number}, \texttt{Cabin\_Side}
\subsubsection{\texttt{Passenger Groups}}
\label{sec:orgfe01226}
A partir de \texttt{PassengerId} foi criada a variável \texttt{Group} permitindo identificar passageiros viajando em conjunto.
\subsection{Transformações Numéricas}
\label{sec:orgb6ad77f}
Foram avaliadas múltiplas transformações \texttt{log1p}, \texttt{sqrt}, \texttt{square} e \texttt{Yeo-Johnson}. As melhores transformações foram selecionadas via validação cruzada estratificada.
\section{Pipeline de Modelagem}
\label{sec:orga8b9b12}
\subsection{Estratégia Geral}
\label{sec:org23b0f3e}
O modelo final utiliza \emph{stacking ensemble} segundo a arquitetura:

\begin{verbatim}
Base Models
    ↓
Probabilidades Out-Of-Fold
    ↓
Meta Learner (Logistic Regression)
    ↓
Predição Final
\end{verbatim}
\section{Encoding Categórico}
\label{sec:orga0af9a7}
As variáveis categóricas foram codificadas utilizando \emph{OrdinalEncoder}
\section{Modelos Utilizados}
\label{sec:org284bdec}
\subsection{CatBoost}
\label{sec:org78aae90}

Modelo baseado em \emph{gradient boosting} com suporte nativo a variáveis categóricas.

Características:
\begin{itemize}
\item excelente desempenho tabular;
\item robustez contra \emph{overfitting};
\item alta capacidade de modelagem não-linear.
\end{itemize}
\subsection{LightGBM}
\label{sec:orge11af59}
\emph{Gradient boosting} baseado em árvores \emph{histogram-based}.

Características:
\begin{itemize}
\item treinamento extremamente rápido;
\item alta escalabilidade;
\item excelente desempenho em \emph{datasets} tabulares.
\end{itemize}
\subsection{XGBoost}
\label{sec:orgd94003b}
\emph{Framework} clássico de \emph{gradient boosting}.

Características:
\begin{itemize}
\item regularização forte;
\item alta estabilidade;
\item ótimo desempenho em competições Kaggle.
\end{itemize}
\subsection{ExtraTrees}
\label{sec:orgce36b7f}
Modelo \emph{ensemble} baseado em árvores extremamente aleatórias.

Características:
\begin{itemize}
\item redução de variância;
\item complementaridade em relação aos modelos boosting.
\end{itemize}
\section{Otimização Bayesiana com Optuna}
\label{sec:orgc3c63f3}
Os hiperparâmetros foram ajustados utilizando \texttt{Optuna} (\emph{Bayesian Optimization} + \emph{Tree-structured Parzen Estimator}, \emph{tree pruning} e busca adaptativa de hiperparâmetros).
\section{Validação Cruzada}
\label{sec:org069de90}

Foi utilizada \emph{Stratified K-Fold Cross Validation} com \texttt{k = 10} para reduzir variância, preservar distribuição da classe alvo e estimar generalização real do modelo.
\section{\emph{Ensemble Stacking}}
\label{sec:org2b0ef93}
O ensemble final foi construído utilizando \emph{StackingClassifier}

\begin{verbatim}
CatBoost
LightGBM
XGBoost
ExtraTrees
    ↓
Logistic Regression
\end{verbatim}

O meta-modelo aprende a combinar as probabilidades dos modelos base.
\section{Otimização de \emph{Threshold}}
\label{sec:orgbf1a893}
Ao invés do threshold padrão \texttt{0.5}, foi realizada busca exaustiva no intervalo \texttt{[0.30, 0.70]} para maximizar acurácia OOF e melhorar calibração das classes.
\section{Métricas Utilizadas}
\label{sec:org78d2562}

\subsection{Accuracy}
\label{sec:org5cd8ca2}
\begin{verbatim}
Accuracy =
Predições corretas / Total
\end{verbatim}

Métrica principal da competição: RMSE.
\subsection{F1 Score}
\label{sec:org8bdb1a2}
Balanceia precisão e recall. Importante em casos de desbalanceamento.
\subsection{ROC AUC}
\label{sec:orga628520}
Mede capacidade discriminativa probabilística.
\section{Resultados}
\label{sec:org2422f5b}
O pipeline desenvolvido tipicamente obtém:

\begin{center}
\begin{tabular}{ll}
Modelo & Acurácia Esperada\\
\hline
CatBoost simples & {[}0.79, 0.81]\\
Ensemble básico & {[}0.81, 0.82]\\
Stack + Optuna & {[}0.82, 0.84]\\
\end{tabular}
\end{center}
\end{document}
